\section{Технический проект}
\subsection{Общая характеристика организации решения задачи}

Необходимо спроектировать и разработать продукционную систему Маркова, которая моделирует применение правил на входной строке и выводит результаты в графическом интерфейсе.  

Система представляет собой набор взаимосвязанных окон, содержащих:
\begin{itemize}
	\item исходный текст;
	\item правила Маркова (загрузка из файла и ручной ввод);
	\item пошаговое отображение изменений текста;
	\item статистику работы алгоритма;
	\item сохранение журнала выполнения шагов в файл.
\end{itemize}

\subsection{Обоснование выбора технологии проектирования}

На современном рынке программных средств для разработки настольных приложений существует множество инструментов. Для поставленной задачи наилучшим выбором является использование языка Python с библиотекой CustomTkinter для реализации GUI и объектно-ориентированного подхода для построения продукционной системы.

\subsubsection{Описание используемых технологий и языков программирования}

\subsubsection{Язык программирования Python}

Python — высокоуровневый интерпретируемый язык программирования, широко используемый для разработки десктопных и веб-приложений, научных вычислений и машинного обучения.  

\paragraph{Достоинства языка Python}
\begin{itemize}
	\item Простота синтаксиса, быстрый прототипинг.
	\item Кроссплатформенность.
	\item Большое количество библиотек для работы с текстом, GUI и данными.
\end{itemize}

\paragraph{Недостатки языка Python}
\begin{itemize}
	\item Интерпретируемый язык, низкая скорость исполнения по сравнению с C/C++.
	\item Ограниченные возможности для создания нативных GUI без сторонних библиотек.
\end{itemize}

\subsubsection{Библиотека CustomTkinter}

CustomTkinter — расширение стандартной библиотеки Tkinter, позволяющее создавать современные GUI-приложения с темной и светлой темой, а также кастомными элементами управления.

\paragraph{Достоинства CustomTkinter}
\begin{itemize}
	\item Современный внешний вид элементов интерфейса.
	\item Простая интеграция с Python.
	\item Поддержка темной и светлой темы.
\end{itemize}

\paragraph{Недостатки CustomTkinter}
\begin{itemize}
	\item Ограниченная документация по сравнению с Tkinter.
	\item Требует установки дополнительного пакета.
\end{itemize}

\subsection{Диаграмма компонентов и схема обмена данными между файлами компонента}

Диаграмма компонентов отражает архитектуру системы:
\begin{itemize}
	\item GUI-приложение (окна, виджеты);
	\item класс \texttt{MarkovRules} (правила Маркова);
	\item класс \texttt{MarkovSystem} (логика применения правил);
	\item файловая система (хранение правил и журнала).
\end{itemize}

\begin{figure}[ht]
	\center{\includegraphics[width=1\linewidth]{components}}
	\caption{Диаграмма компонентов продукционной системы Маркова}
	\label{comp:image}
\end{figure}

Схема обмена данными между компонентами представлена на рисунке \ref{data:image}. Пользователь через GUI передаёт исходный текст и правила классу \texttt{MarkovSystem}, который применяет правила, обновляет текст, ведёт журнал и формирует статистику. Результаты отображаются в интерфейсе и сохраняются в файл.

\begin{figure}[ht]
	\center{\includegraphics[width=1\linewidth]{dataflow}}
	\caption{Схема обмена данными между компонентами}
	\label{data:image}
\end{figure}

\subsection{Диаграмма размещения}

Диаграмма размещения показывает физическую структуру системы: пользовательский компьютер с GUI-приложением и файловую систему для хранения правил и журналов.

\begin{figure}[ht]
	\center{\includegraphics[width=0.6\linewidth]{deployment}}
	\caption{Диаграмма размещения компонентов системы}
	\label{place:image}
\end{figure}

\subsection{Содержание информационных блоков. Основные сущности}

Основные сущности программы:
\begin{itemize}
	\item {Пользователь}
	\item {Правило Маркова}
	\item {Система Маркова}
	\item {Журнал шагов}
	\item {Результат выполнения}
\end{itemize}


\paragraph{Описание сущностей и их атрибутов}

\renewcommand{\arraystretch}{0.8}
\begin{xltabular}{\textwidth}{|X|p{3cm}|X|X|}
	\caption{Атрибуты сущности \texttt{Пользователь}\label{user:table}}\\
	\hline \centrow Поле & \centrow Тип & \centrow Обязательное & \centrow Описание \\ \hline
	\endfirsthead
	\caption*{Продолжение таблицы \ref{user:table}}\\
	\hline \centrow Поле & \centrow Тип & \centrow Обязательное & \centrow Описание \\ \hline
	\finishhead
	name & String & true & Имя пользователя \\ \hline
\end{xltabular}
\renewcommand{\arraystretch}{1.0}

\renewcommand{\arraystretch}{0.8}
\begin{xltabular}{\textwidth}{|X|p{3cm}|X|X|}
	\caption{Атрибуты сущности \texttt{Правило Маркова}\label{rule:table}}\\
	\hline \centrow Поле & \centrow Тип & \centrow Обязательное & \centrow Описание \\ \hline
	\endfirsthead
	\caption*{Продолжение таблицы \ref{rule:table}}\\
	\hline \centrow Поле & \centrow Тип & \centrow Обязательное & \centrow Описание \\ \hline
	\finishhead
	left & String & true & Левая часть правила \\ \hline
	right & String & true & Правая часть правила \\ \hline
	stop & Bool & false & Признак завершения работы системы \\ \hline
\end{xltabular}
\renewcommand{\arraystretch}{1.0}

\renewcommand{\arraystretch}{0.8}
\begin{xltabular}{\textwidth}{|X|p{3cm}|X|X|}
	\caption{Атрибуты сущности \texttt{Система Маркова}\label{system:table}}\\
	\hline \centrow Поле & \centrow Тип & \centrow Обязательное & \centrow Описание \\ \hline
	\endfirsthead
	\caption*{Продолжение таблицы \ref{system:table}}\\
	\hline \centrow Поле & \centrow Тип & \centrow Обязательное & \centrow Описание \\ \hline
	\finishhead
	rules & List[MarkovRules] & true & Список правил \\ \hline
	steps & Integer & true & Количество выполненных шагов \\ \hline
	applied\_rules & Dict & true & Статистика применения правил \\ \hline
	lengths & List[Int] & true & Длина строки на каждом шаге \\ \hline
\end{xltabular}
\renewcommand{\arraystretch}{1.0}

\renewcommand{\arraystretch}{0.8}
\begin{xltabular}{\textwidth}{|X|p{3cm}|X|X|}
	\caption{Атрибуты сущности \texttt{Журнал шагов}\label{log:table}}\\
	\hline \centrow Поле & \centrow Тип & \centrow Обязательное & \centrow Описание \\ \hline
	\endfirsthead
	\caption*{Продолжение таблицы \ref{log:table}}\\
	\hline \centrow Поле & \centrow Тип & \centrow Обязательное & \centrow Описание \\ \hline
	\finishhead
	log & List[String] & true & Хронология изменений текста \\ \hline
\end{xltabular}
\renewcommand{\arraystretch}{1.0}

\renewcommand{\arraystretch}{0.8}
\begin{xltabular}{\textwidth}{|X|p{3cm}|X|X|}
	\caption{Атрибуты сущности \texttt{Результат выполнения}\label{result:table}}\\
	\hline \centrow Поле & \centrow Тип & \centrow Обязательное & \centrow Описание \\ \hline
	\endfirsthead
	\caption*{Продолжение таблицы \ref{result:table}}\\
	\hline \centrow Поле & \centrow Тип & \centrow Обязательное & \centrow Описание \\ \hline
	\finishhead
	final\_text & String & true & Текст после применения всех правил \\ \hline
	steps & Integer & true & Общее количество шагов \\ \hline
	most\_common\_rule & String & false & Наиболее часто применяемое правило \\ \hline
	avg\_length & Float & false & Средняя длина текста \\ \hline
\end{xltabular}
\renewcommand{\arraystretch}{1.0}